\documentclass[11pt]{amsart}
\usepackage[margin=1in]{geometry}
\usepackage{amsmath,amssymb,amsthm,mathtools,booktabs,array}
\usepackage[T1]{fontenc}
\usepackage{lmodern,microtype}
\usepackage{xurl}
\usepackage[hidelinks]{hyperref}
\hypersetup{pdftitle={Real and simple zeros of the squared Eulerian triangle},pdfauthor={Claude Opus 5.5 and GPT-6 Astra}}
\newtheorem{theorem}{Theorem}[section]
\newtheorem{lemma}[theorem]{Lemma}
\newtheorem{proposition}[theorem]{Proposition}
\newtheorem{corollary}[theorem]{Corollary}
\theoremstyle{definition}
\newtheorem{definition}[theorem]{Definition}
\newtheorem{remark}[theorem]{Remark}
\newcommand{\RR}{\mathbb R}
\newcommand{\NN}{\mathbb N}
\newcommand{\E}{\mathbb E}
\DeclareMathOperator{\Var}{Var}
\DeclareMathOperator{\sgn}{sgn}
\newcommand{\eps}{\varepsilon}
\title[The squared Eulerian triangle]{Real and simple zeros of the squared Eulerian triangle}
\author{Claude Opus 5.5 \and GPT-6 Astra}
\date{1 November 2026}
\begin{document}
\begin{abstract}
Let $A(n,k)$ be the Eulerian numbers counting permutations with $k$ ascending runs, and let $B=A\cdot A$ be the ordinary matrix square of the Eulerian triangle. We prove that every polynomial $P_n(x)=\sum_k B(n,k)x^k$, for $n\ge1$, has a simple zero at zero and $n-1$ distinct negative zeros. This resolves the Eulerian part of Mao and Wang's Conjecture~4.1. The proof combines a monotone phase in an exact pole expansion, uniform bounds for the remaining poles, and sign changes in the two extreme ranges. In the transition at large negative arguments, an exact two-dimensional weight comparison separates close pairs of zeros. Integer sign certificates cover $n\le250$; the remaining verification consists of explicit rational inequalities for universal constants.
\end{abstract}
\maketitle

\section{Introduction}
Define $A(0,0)=1$, put $A(n,k)=0$ outside $0\le k\le n$, and set
\begin{equation}\label{eq:eulerian-recurrence}
A(n,k)=(n-k+1)A(n-1,k-1)+kA(n-1,k)\qquad(n\ge1).
\end{equation}
For $n\ge1$, $A(n,k)$ counts permutations of $n$ letters with $k-1$ descents, or equivalently with $k$ ascending runs. Write
\[
A_n(x)=\sum_k A(n,k)x^k,\qquad
B(n,k)=\sum_{j=k}^n A(n,j)A(j,k),\qquad
P_n(x)=\sum_k B(n,k)x^k.
\]
The square defining $B$ is ordinary matrix multiplication. Mao and Wang conjectured that the row polynomials of this square have only real nonpositive zeros \cite[Conjecture~4.1]{MaoWang2026}. Their conjecture also treats a separate Delannoy triangle. We prove its Eulerian part, including simplicity.

\begin{theorem}\label{thm:main}
For every integer $n\ge1$, $P_n$ has exactly $n$ distinct real zeros: one at zero and $n-1$ in $(-\infty,0)$. The polynomial $P_0=1$ has no zeros.
\end{theorem}

For example,
\[
\begin{aligned}
P_1(x)&=x, &P_2(x)&=2x+x^2,\\
P_3(x)&=6x+8x^2+x^3, &
P_4(x)&=24x+66x^2+22x^3+x^4.
\end{aligned}
\]
The Eulerian transformation $T(t^j)=A_j(x)$ gives $P_n=T(A_n)$. General real-rootedness preservation by $T$ is false \cite{BrandenJochemko2022}; the positive cone on which preservation is known \cite{Athanasiadis2025} does not contain $A_n$ for $n\ge3$. Our proof uses the particular Eulerian coefficients.

The zeros occupy very different scales. We use the coordinate $x=-e^{-s}$. A pole expansion has a strictly increasing phase, and its nearest conjugate poles determine the sign for $|s|\le4\sqrt n$. Explicit coefficient estimates handle the small negative zeros. At the other end, a four-term comparison gives pairs of large negative zeros. Between this comparison and the pole range, the maximizing terms of an exact double sum move along an ordered lattice path. Single transitions give one sign change. Two nearby transitions give two sign changes because their mixed coefficient is strictly smaller than the product of the two independent coefficients. This last strict inequality persists after every remaining term is included.

The argument for $n\ge250$ is uniform. A finite family of rational sign certificates supplies the initial range. Every numerical constant used in the estimates is an explicit rational bound, checked by the standard-library supplement described in Appendix~\ref{sec:verification}. Approximations of polynomial zeros are not used in the uniform argument.

\section{Eulerian identities and a monotone phase}

\begin{lemma}\label{lem:eulerian}
For $m\ge1$, $A_m$ has a simple zero at zero and $m-1$ simple negative zeros. It satisfies
\[
A_m(x)=x^{m+1}A_m(1/x),\qquad
\frac{A_m(x)}{(1-x)^{m+1}}=\sum_{h\ge1}h^m x^h\quad(|x|<1).
\]
\end{lemma}
\begin{proof}
The recurrence gives
\[
A_m=x\{(1-x)A_{m-1}'+mA_{m-1}\}.
\]
At each negative zero of $A_{m-1}$ the expression in braces has the sign of $A_{m-1}'$. These signs alternate. Its value at zero is positive, and its leading coefficient is positive. The intermediate value theorem, including the leftmost interval and the interval ending at zero, gives all $m-1$ negative zeros and proves their simplicity by degree. The initial case is $A_1=x$. Reversing a permutation proves the reciprocal identity. Applying $x\,d/dx$ to the geometric series and using the same recurrence proves the series identity. See also \cite{Petersen2015}.
\end{proof}

\begin{lemma}\label{lem:bonferroni}
For positive integers $m,k$,
\begin{equation}\label{eq:bonferroni}
k^m-(m+1)(k-1)^m\le A(m,k)\le k^m.
\end{equation}
The sequence $A(m,k)^{1/m}$, with its zero endpoints, is concave in $k$.
\end{lemma}
\begin{proof}
The density at $k$ of a sum of $m+1$ independent uniforms on $[0,1]$ is $A(m,k)/m!$. Equivalently, the $m$-dimensional volume of the section of $[0,1]^{m+1}$ by $\sum x_i=k$ is $\sqrt{m+1}\,A(m,k)/m!$. In the positive orthant without the upper coordinate bounds, the corresponding section has volume $\sqrt{m+1}\,k^m/m!$. Each of the $m+1$ excluded sets $x_i>1$ has volume $\sqrt{m+1}\,(k-1)^m/m!$. The upper bound and the first union bound prove \eqref{eq:bonferroni}. Brunn's concavity theorem for parallel sections of a convex body gives the final assertion \cite{Gardner2002}.
\end{proof}

Define the factorial-weighted polynomial
\[
F_n(t)=\sum_{j=1}^n j!A(n,j)t^j\quad(n\ge1),\qquad F_0=1.
\]
\begin{lemma}\label{lem:factorial}
For $n\ge1$, $F_n$ has $n$ simple nonpositive zeros, exactly one of which is zero. Thus
\begin{equation}\label{eq:factorial-factorization}
F_n(t)=n!t\prod_{i=1}^{n-1}(t+\beta_i),\qquad\beta_i>0.
\end{equation}
\end{lemma}
\begin{proof}
Writing $f=F_{n-1}$ and $\theta=t\,d/dt$, coefficient comparison in \eqref{eq:eulerian-recurrence} gives
\[
F_n=t\{f'+(1+\theta)(n-\theta)f\}
 =t\{nf+[1+(n-2)t]f'-t^2f''\}.
\]
The cases $n=1,2$ are $t$ and $t+2t^2$. Assume the result at $n-1$. At any critical point $c$ of $f$, the expression in braces is $nf(c)-c^2f''(c)$, which has the sign of $f(c)$: each critical point lies between consecutive simple real zeros, and is a positive maximum or a negative minimum. There are $n-2$ such points, their signs alternate, and the last sign is negative. The expression in braces has degree $n-1$, positive leading coefficient, and value $f'(0)=1$ at zero. Its signs at the critical points, at negative infinity, and at zero give $n-1$ distinct negative zeros. Multiplication by $t$ completes the induction.
\end{proof}

For real $s$ put
\[
C=1+e^{-s},\qquad z_q=s+iq\pi\quad(q\text{ odd}),\qquad
R_s(z)=\frac{F_n(C/z)}{z}.
\]
\begin{lemma}\label{lem:poles}
For $n\ge1$,
\begin{equation}\label{eq:pole-expansion}
P_n(-e^{-s})=C\sum_{q\text{ odd}}R_s(z_q).
\end{equation}
The series is absolutely convergent and its terms pair by conjugation.
\end{lemma}
\begin{proof}
The symmetrically paired partial-fraction identity
\[
\sum_{q\text{ odd}}\frac1{s+iq\pi}=\frac12\tanh(s/2)
 =\frac1{1+e^{-s}}-\frac12
\]
and its derivatives give, for $j\ge1$,
\[
A_j(-e^{-s})=j!C^{j+1}\sum_{q\text{ odd}}z_q^{-j-1}.
\]
One may first differentiate for $s>0$ and use Lemma~\ref{lem:eulerian}, then continue along the real line. Multiply by $A(n,j)$ and sum over $j$. Since $F_n(0)=0$, the resulting terms are $O(q^{-2})$, which proves absolute convergence.
\end{proof}

\begin{lemma}\label{lem:phase}
The continuous argument of the nearest upper-pole term is
\begin{equation}\label{eq:phase}
\phi(s)=-(n+1)\arg(s+i\pi)
 +\sum_{i=1}^{n-1}\arg(s+C/\beta_i+i\pi),
\end{equation}
where every argument on the right lies in $(0,\pi)$. It is strictly increasing from $-(n+1)\pi$ to zero. In particular every level $j\pi$, $-n\le j\le-1$, is crossed exactly once.
\end{lemma}
\begin{proof}
Factorization \eqref{eq:factorial-factorization} gives the displayed argument. Put $a=C/\beta_i$ and $\lambda=1/(1+e^s)$. The derivative of $s+a$ is $1-\lambda a$, and
\[
\frac1{s^2+\pi^2}-\frac{1-\lambda a}{(s+a)^2+\pi^2}
=\frac{a\{a+2s+\lambda(s^2+\pi^2)\}}
 {(s^2+\pi^2)((s+a)^2+\pi^2)}>0.
\]
For $s\ge0$ positivity is immediate. For $s<0$, use $\lambda\ge1/2$ and
\[
2s+\lambda(s^2+\pi^2)\ge\tfrac12\{(s+2)^2+\pi^2-4\}>0.
\]
Consequently $\phi'(s)\ge2\pi/(s^2+\pi^2)>0$. The two limits follow directly from \eqref{eq:phase}.
\end{proof}

The other exact representation used below follows by reversing the coefficients in the matrix product:
\begin{equation}\label{eq:double-sum}
\frac{P_n(-X)}{(-X)^n}
=\sum_{\substack{k,h\ge1\\k+h\le n+1}}
 (-1)^{k+h-2}A(n,k)A(n-k+1,h)X^{-(k+h-2)}\quad(X>0).
\end{equation}
Indeed the coefficient of $X^{-r}$ before its sign is
\[
B(n,n-r)=\sum_{k+h=r+2}A(n,k)A(n-k+1,h),
\]
by Eulerian symmetry. Its positive comparison weight will be
\begin{equation}\label{eq:ideal-weight}
W_u(k,h)=k^n h^{n-k+1}e^{-(k+h-2)u},\qquad X=e^u.
\end{equation}

% Uniform phase signs for the Eulerian-square polynomial.
\section{Uniform phase signs in the central region}\label{sec:bulk}

Throughout this section $n\ge250$. We use the real-rootedness of the
factorial-weighted polynomial
\[
 F_n(t)=\sum_{j=1}^n j!A(n,j)t^j
       =n!t\prod_{i=1}^{n-1}(t+\beta_i),\qquad \beta_i>0.
\]
For real $s$, set $C_s=1+e^{-s}$ and
\[
 \mathcal R_{n,s}(z)=\frac{F_n(C_s/z)}z,\qquad
 z_q=s+iq\pi\quad(q\text{ odd}).
\]
The pole expansion and its continuous nearest-pole phase are
\[
 P_n(-e^{-s})=C_s\sum_{q\text{ odd}}\mathcal R_{n,s}(z_q),
 \qquad
 \phi(s)=\arg\mathcal R_{n,s}(s+i\pi).
\]
The sum is absolutely convergent. All phase comparisons below use the
continuous argument specified by the factorization of $F_n$.

\begin{lemma}[Central phase signs]\label{lem:bulk}
For every $n\ge250$ and $|s|\le4\sqrt n$,
\[
 \sum_{q=3,5,7,\ldots}
 \frac{|\mathcal R_{n,s}(s+iq\pi)|}
      {|\mathcal R_{n,s}(s+i\pi)|}<1.
\]
In particular, if $\phi(s)=j\pi$ in this interval, then
\[
 \operatorname{sgn}P_n(-e^{-s})=(-1)^j.
\]
One may use the upper bounds $.85$ for $0\le s\le3\sqrt n$,
$.33$ for $3\sqrt n\le s\le4\sqrt n$, and $.88$ for
$-4\sqrt n\le s\le0$.
\end{lemma}

We prove the three estimates separately. The finite arithmetic checks
used in the proof involve only rational constants; their precise
inequalities and bounding series are given at the end of the section.

\subsection{A weighted mean and the positive half of the interval}

\begin{lemma}\label{lem:weightedmean}
For $n\ge2$ and $t\ge e/n$,
\[
 \frac{tF_n'(t)}{F_n(t)}\ge\frac{n+1}{2}.
\]
\end{lemma}
\begin{proof}
Pair the terms indexed by $j$ and $\ell$ with $j+\ell=n+1$ and
$j>\ell$. Eulerian symmetry gives $A(n,j)=A(n,\ell)$. The geometric
mean of $\ell+1,\ldots,j$ is at least $(n!)^{1/(n-1)}$: these intervals
have the same center, and deleting an outer symmetric pair increases
the average logarithm, since the products of the symmetric pairs
increase toward the center. Moreover,
\[
 \log(n!)\ge\int_1^n\log x\,dx
 =n\log n-n+1\ge(n-1)(\log n-1).
\]
Thus $j!t^j\ge\ell!t^\ell$. Each pair makes a nonnegative contribution
to $\sum_j(j-(n+1)/2)j!A(n,j)t^j$, proving the assertion.
\end{proof}

Fix $s\ge0$, put $C=C_s$ and $a_i=C/\beta_i$, and write
\[
 R_q=\frac{|\mathcal R_{n,s}(s+iq\pi)|}
 {|\mathcal R_{n,s}(s+i\pi)|}.
\]
Direct factorization gives
\begin{equation}\label{eq:bulk-positive-ratio}
 R_q=\left(\frac{s^2+\pi^2}{s^2+q^2\pi^2}\right)^{(n+1)/2}
 \prod_{i=1}^{n-1}
 \left(\frac{(s+a_i)^2+q^2\pi^2}{(s+a_i)^2+\pi^2}\right)^{1/2}.
\end{equation}
For $0<s\le3\sqrt n$, $C/s\ge e/n$, so Lemma~\ref{lem:weightedmean}
implies
\[
 \sum_i\frac{a_i}{s+a_i}\ge\frac{n-1}{2}.
\]
At least $(n-1)/3$ of the $a_i$ are at least $s/3$: if $M$ have this
property, the left side is at most $M+(n-1-M)/4$.
Set
\[
 v=\frac{s^2}{\pi^2},\quad c=\frac{16}{9},\quad
 M_0=\frac{n-1}{3},\quad V=\frac{9n}{\pi^2}.
\]
Bounding the $M$ distant factors separately in
\eqref{eq:bulk-positive-ratio} gives
\begin{equation}\label{eq:positive-fq}
 R_q\le f_q(v):=\frac{v+1}{v+q^2}
 \left\{\frac{(v+1)(cv+q^2)}{(v+q^2)(cv+1)}\right\}^{M_0/2}.
\end{equation}
The sign of the logarithmic derivative of $f_q$ is the sign of
\[
 \left[c^2+\frac{M_0}{2}c(c-1)\right]v^2
 +c(1+q^2)v+q^2\left[1-\frac{M_0}{2}(c-1)\right].
\]
Its first two coefficients are positive. Hence $f_q$ decreases and
then increases, or increases throughout, and its maximum on $[0,V]$
is attained at an endpoint. At zero, $f_q(0)=q^{-2}$. Using
$\log(1-x)\le-x$ at the other endpoint gives
\[
 f_q(V)\le e^{-E_n(q)},\qquad
 E_n(q)=\frac{7n(n-1)(q^2-1)\pi^2}
 {6(9n+q^2\pi^2)(16n+\pi^2)}.
\]
For fixed $q$, $E_n(q)$ increases with $n$. At $n=250$,
\[
 E_n(3)>.60,\quad E_n(5)>1.70,\quad
 E_n(7)>3.11,\quad E_n(9)>4.60.
\]
Thus the first three ratios are less than $.55,.19,.045$, respectively,
and $f_9(V)<1/81$. Differentiation in $y=q^2$ shows that
$yf_{\sqrt y}(V)$ decreases whenever
\[
 y\ge\frac{cV}{M_0(c-1)/2-1}<13.
\]
Consequently $f_q(V)\le q^{-2}$ for $q\ge9$, and the endpoint argument
extends this inequality to every $v\in[0,V]$. Since
\[
 \sum_{q=9,11,\ldots}q^{-2}
 =\frac{\pi^2}{8}-1-\frac19-\frac1{25}-\frac1{49}<.063,
\]
we obtain $\sum_{q\ge3,\ q\text{ odd}}R_q<.848<.85$. At $s=0$ the
same result follows directly from $R_q\le q^{-2}$.

For the remaining positive interval, $3\sqrt n\le s\le4\sqrt n$,
Lemma~\ref{lem:weightedmean} still applies because
$C/s\ge1/(4\sqrt n)\ge e/n$. Put
\[
 p_i=\frac{a_i}{s+a_i},\quad t_i=1-p_i,\quad
 A=\frac{q^2\pi^2}{s^2},\quad B=\frac{\pi^2}{s^2},
\]
so $\sum_i p_i\ge(n-1)/2$. Define
\[
 h_q(t)=\frac12\log\frac{1+At^2}{1+Bt^2},\qquad D_q=h_q(1).
\]
If $A\le3$ and $B\le.005$, then $h_q(t)\le tD_q$ on $[0,1]$.
Indeed, for
$\psi(v)=v/(1+v)-\tfrac12\log(1+v)$,
\[
 t^2\left(\frac{h_q(t)}t\right)'=\psi(At^2)-\psi(Bt^2).
\]
The function $\psi$ increases on $[0,1]$ and decreases on $[1,3]$.
If $At^2\le1$, the difference is nonnegative. Otherwise
$\psi(At^2)\ge\psi(3)>.05$, whereas $\psi(Bt^2)\le B/2\le.0025$.
It follows from \eqref{eq:bulk-positive-ratio} that
\begin{equation}\label{eq:positive-strong}
 \log R_q\le-\frac{n+3}{2}D_q
 \qquad\left(\frac{q\pi}{s}\le\sqrt3\right).
\end{equation}
Let $Q$ be the largest odd integer at most $\sqrt3s/\pi$. Then
$Q\ge25$ and
\[
 \frac{Q\pi}{s}>\sqrt3\frac{Q}{Q+2}
 \ge\sqrt3\frac{25}{27}>1.6.
\]
The inequality $\log(1+x)\ge2x/(2+x)$ in
\eqref{eq:positive-strong} gives, for $q=3,5,7,9$,
\[
 -\log R_q\ge
 \frac{(q^2-1)9.86}{64+2(q^2+1)9.87/250}.
\]
The corresponding four upper bounds sum to less than $.326$.
For $11\le q\le Q$ the same inequality gives
\[
 -\log R_q\ge\alpha(q^2-1),\qquad \alpha=\frac{9.86}{160.16},
\]
and hence
\[
 \sum_{11\le q\le Q,\ q\text{ odd}}R_q
 \le\frac{e^{-120\alpha}}{1-e^{-48\alpha}}<.001.
\]
For later use, comparison of the individual root factors yields, for
any real $s$, any odd $Q$, and $q\ge Q$,
\begin{equation}\label{eq:universal-pole-tail}
 \frac{R_q}{R_Q}\le\left(\frac Qq\right)^2
 \exp\!\left(\frac{(n+1)s^2}{2Q^2\pi^2}\right),\qquad
 \sum_{q>Q,\ q\text{ odd}}R_q
 \le\frac Q2R_Q\exp\!\left(\frac{(n+1)s^2}{2Q^2\pi^2}\right).
\end{equation}
Here the last inequality uses
$\sum_{q>Q,\ q\text{ odd}}q^{-2}\le1/(2Q)$.
At the present $Q$,
$D_Q\ge\tfrac12\log(3.56/1.005)>.62$, so $R_Q<e^{-.31n}$.
The extra exponential in \eqref{eq:universal-pole-tail} is at most
$e^{.197n}$ and $Q/2<2\sqrt n$. The resulting tail is less than $.001$
for $n\ge250$. This proves the positive-interval bound $.33$.

\subsection{A uniform variance estimate}

We next record the distributional estimate needed on the negative
half of the interval. Write $(n)_\ell=n!/(n-\ell)!$.

\begin{lemma}\label{lem:variance}
Let $n\ge250$, $\pi\le u\le4\sqrt n$, and
\[
 B=\frac{1+e^u}{u},\qquad Y=\log(nB).
\]
If $L$ has probabilities
\[
 p_\ell\ \propto\ \frac{A(n,\ell+1)}{(n)_\ell B^\ell},
 \qquad 0\le\ell\le n-1,
\]
then $L$ is a sum of independent Bernoulli variables and
\[
 \operatorname{Var}L<\frac{8n}{Y^2}.
\]
\end{lemma}
\begin{proof}
The real-rootedness of $F_n$, coefficient reversal, and positive
tilting show that $L$ is a Bernoulli sum. More explicitly, its success
probabilities are $\beta_i/(B+\beta_i)$. In particular its probabilities
are log-concave. Let $m$ be a mode and put $N=n-m$.

The sequence $u_k=A(n,k)^{1/n}$, with zero endpoints, is concave.
Indeed, the $n$-dimensional volume of the section
$\sum_{i=1}^{n+1}x_i=k$ of $[0,1]^{n+1}$ is
$\sqrt{n+1}A(n,k)/n!$; Brunn's concavity theorem applies
\cite{Gardner2002}. Thus $u_{k+1}/u_k\le(k+1)/k$.
Eulerian unimodality places $m$ before the midpoint. If $m\ge1$,
\[
 B(n-m+1)\le\frac{A(n,m+1)}{A(n,m)}
 \le(1+1/m)^n\le e^{n/m}.
\]
Consequently
\begin{equation}\label{eq:bulk-mode-mean}
 m\le\frac n{Y-\log2},\qquad
 \mu:=\mathbb EL\le\frac n{Y-\log2}+1,
\end{equation}
where the second bound is Darroch's mode theorem
\cite{Darroch1964}. The first bound also holds when $m=0$.

Put $h=Y^2/n$. Elementary endpoint estimates give
\begin{equation}\label{eq:variance-parameters}
 Y>7.5,\quad Y/n<.26,\quad h<17,\quad
 N/n>.852,\quad\log(N/n)>-.16.
\end{equation}
For the two upper bounds, $Y$ increases with $u$ and at $u=4\sqrt n$,
\[
 Y=4\sqrt n+\tfrac12\log n-\log4+
       \log(1+e^{-4\sqrt n}).
\]
The resulting bounds on $Y/n$ and $Y/\sqrt n$ decrease after $n=250$;
the endpoint gives $Y<65$ and $Y/\sqrt n<\sqrt{17}$.
The lower bound follows at $n=250,u=\pi$.

Set
\[
 a=\frac{p_{m-1}}{p_m},\quad b=\frac{p_{m+1}}{p_m},\quad
 c=(BN)^{1/n}-1,\quad e=1-[B(N+1)]^{-1/n},\quad
 J=\lfloor n/Y\rfloor\ge3,
\]
with missing adjacent probabilities set equal to zero. Concavity of
$u_k$ gives the tangent estimates
\begin{align}
 \frac{p_{m+j}}{p_m}&\le b\,e^{-hj(j-1)/7}
       &&(1\le j\le J),\label{eq:variance-forward}\\
 \frac{p_{m-j}}{p_m}&\le a\,e^{-hj(j-1)/3}
       &&(1\le j\le m),\label{eq:variance-backward}\\
 ab&\le e^{-h/2}.\label{eq:variance-adjacent}
\end{align}
Here are the details retaining the adjacent factors. If
$c_0=u_{m+2}/u_{m+1}-1\le c$, concavity implies
$u_{m+j+1}/u_{m+1}\le1+jc_0$. After subtracting $\log b$, the
$c_0$-dependent expression
$n[\log(1+jc_0)-\log(1+c_0)]$ is increasing in $c_0$.
Therefore
\[
 \log\frac{p_{m+j}}{bp_m}
 \le n[\log(1+jc)-j\log(1+c)]
      +\sum_{i=0}^{j-1}-\log(1-i/N)
 \le-\left\{\frac{nc^2}{(1+Jc)(1+c)}-\frac1{N-J}\right\}
       \frac{j(j-1)}2.
\]
The elementary bounds
\[
 c\ge\left(1-\frac{.16}{7.5}\right)\frac Yn,
 \quad c<.298,\quad Jc<1.15,\quad (N-J)/n>.719
\]
make the bracket exceed $2h/7$.
For the backward direction the actual decrement
$d=1-u_m/u_{m+1}$ is at least $e$. Subtracting $\log a$ leaves
$n[\log(1-jd)-\log(1-d)]$, which decreases in $d$. Hence
\[
 \log\frac{p_{m-j}}{ap_m}
 \le n[\log(1-je)-j\log(1-e)]
       +\sum_{i=0}^{j-1}\log(1+i/(N+1))
 \le-\left(ne^2-\frac1{N+1}\right)\frac{j(j-1)}2.
\]
Since $1-e^{-t}\ge.88t$ for $0\le t\le.26$, the bracket exceeds
$2h/3$. This argument is needed only for $m\ge1$, so $N+1\le n$.
Finally, concavity gives
$u_mu_{m+2}/u_{m+1}^2\le(1-d)(1+d)\le1-e^2$, and thus
\[
 ab\le(1-e^2)^n\frac{N+1}{N}
 \le\exp\!\left(-ne^2+\log\frac{N+1}{N}\right)<e^{-h/2}.
\]
When $m=0$, the backward assertion is empty and $ab=0$.

The one-step version of the forward estimate gives
$p_{m+J}/p_{m+J-1}\le e^{-2h(J-1)/7}$.
All later ratios are no larger, by log-concavity. For $t=n/Y\ge3$,
\[
 \frac{J(J-1)}{t^2}\ge\frac38,\qquad
 \frac{J-1}{t}\ge\frac12.
\]
The initial amplitude is at most $e^{-3n/56}$ and the subsequent
ratio is less than $e^{-Y/7}<.343$. Summing the geometric tail gives
\begin{equation}\label{eq:variance-far-tail}
 \sum_{j>J}j^2\frac{p_{m+j}}{p_m}
 \le1.3J^2e^{-3n/56}
 \le1.3(n/7.5)^2e^{-3n/56}<.003.
\end{equation}

Let $v=\operatorname{Var}L$. Fourier inversion for a Bernoulli sum gives
\[
 \max_k p_k\le e^{-v}I_0(v),\qquad
 I_0(v)=\sum_{k\ge0}\frac{(v^2/4)^k}{(k!)^2}.
\]
Indeed, the squared modulus of its characteristic function is at most
$\exp[-4v\sin^2(t/2)]$. We shall use
\begin{equation}\label{eq:variance-mass}
 \sqrt v\,e^{-v}I_0(v)<.47\qquad(v\ge.7).
\end{equation}
For $v\ge1$, put $f(v)=\sqrt v\,e^{-v}I_0(v)$. The Bessel equation gives
$f''+2f'+f/(4v^2)=0$. Writing $I_1=I_0'$, the defining series imply
$2I_1(1)<I_0(1)$, so $f'(1)<0$; at every zero of $f'$, one has
$f''<0$. Thus $f$ decreases on $[1,\infty)$, and $f(1)<.47$.
On $[.7,1]$, inequality~\eqref{eq:variance-mass} is certified on
300 rational intervals of length $.001$ by the bounds below.

Define
\[
 S(t)=\sum_{j\ge1}j^2e^{-tj(j-1)},\qquad
 T(t)=\sum_{j\ge2}(j-1)^2e^{-tj(j-1)}.
\]
Equations~\eqref{eq:variance-forward}--\eqref{eq:variance-far-tail}
imply
\[
 v\le\mathbb E(L-m)^2
 \le p_m\{bS(h/7)+aS(h/3)+.003\}
 \le p_m\{S(h/7)+e^{-h/2}S(h/3)+.003\}.
\]
The last inequality maximizes a linear function over
$0\le a,b\le1$, $ab\le e^{-h/2}$, using $S(h/7)\ge S(h/3)$.
If $h\le11$ and $v\ge8/h$, then $v>.7$ and
\[
 v^{3/2}<.47\{S(h/7)+e^{-h/2}S(h/3)+.003\}
          <(8/h)^{3/2},
\]
a contradiction. For $0<h\le2$, the last comparison follows from
$S(t)\le2t^{-3/2}$ for $0<t\le2/3$. To verify this elementary bound,
apply the integral-plus-maximum bound for a positive unimodal sum to
$x^2e^{-tx(x-1)}$. After multiplication by $t^{3/2}$, its integral
increases with $t$ and is less than $1.21$ at $t=2/3$; its maximum
also increases with $t$ and is less than $.75$ there. For the remaining
values of $h\le11$, the six rational interval checks in
Table~\ref{tab:variance-intervals} prove the comparison.

For $11\le h<17$, clamp $L$ to $[m,m+1]$ if $b\ge a$, and to
$[m-1,m]$ otherwise. Write the resulting variable as $Z$, and
$W=L-Z$. In the first case $a\le e^{-h/4}$, and in the second
$b\le e^{-h/4}$. Since $\operatorname{Var}Z\le1/4$,
\[
 v\le\frac14+\mathbb EW^2+\sqrt{\mathbb EW^2}
 \le\frac14+\Delta(h)+\sqrt{\Delta(h)},
\]
where
\[
 \Delta(h)=\max\{T(h/7)+e^{-h/4}S(h/3),\,
                    T(h/3)+e^{-h/4}S(h/7)\}+.003.
\]
This function decreases. On each half-unit interval $[a,b]$ in
$[11,17]$, exact rational bounds give
\[
 8/b-1/4-\Delta(a)>0,\qquad
 \Delta(a)<[8/b-1/4-\Delta(a)]^2.
\]
Consequently $v<8/h$, completing the proof.
\end{proof}

\subsection{Kernels for the negative half of the interval}

Write $s=-u$, $u\ge0$. For $u\ge\pi$ set
\[
 B=\frac{1+e^u}{u},\quad Y=\log(nB),\quad
 r_i=\beta_i/B,\quad p_i=\frac{r_i}{1+r_i},\quad
 w_i=p_i(1-p_i).
\]
Then $\mu=\sum p_i$ and $v=\sum w_i$ are the mean and variance in
Lemma~\ref{lem:variance}. For $q\ge1$ let
\[
 D_q=\frac12\log\frac{u^2+q^2\pi^2}{u^2+\pi^2},\qquad
 g_q(r)=\frac12\log
 \frac{u^2(1-1/r)^2+q^2\pi^2}{u^2(1-1/r)^2+\pi^2}.
\]
The exact pole ratio satisfies
\begin{equation}\label{eq:negative-exact-ratio}
 \log R_q=-(n+1)D_q+\sum_i g_q(r_i).
\end{equation}
For $p=r/(1+r)$ and $w=p(1-p)$, the needed kernel bounds are
\begin{align}
 g_3(r)&\le D_3p+4\log3\,w,\label{eq:kernel-three}\\
 g_q(r)&\le D_qp+[4\log q+(q-1)/2]w,\label{eq:kernel-variance}\\
 g_q(r)&\le\max(q,4\log q)\,p.\label{eq:kernel-mean}
\end{align}
We give proofs to fix the constants.

For $r\ge1$, put $b=1/r\in[0,1]$ and $L=\log q$. Convexity of
$z\mapsto\tfrac12\log[(u^2z+q^2\pi^2)/(u^2z+\pi^2)]$ gives
\[
 g_q(r)\le D_q+(2b-b^2)(L-D_q)
 \le\frac{D_q}{1+b}+\frac{4Lb}{(1+b)^2}.
\]
The last comparison is affine in $D_q\in[0,L]$; at its endpoints it
reduces to $(1-b)^2(b+2)\ge0$ and $b(3-b)\ge0$.
This proves the right half of \eqref{eq:kernel-three} and
\eqref{eq:kernel-variance}.
For $r\le1$, $g_3$ decreases as $u$ increases, so it suffices to take
$u=\pi$. Put
\[
 N_r=1-2r+10r^2,\quad Q_r=1-2r+2r^2,\quad
 H(r)=\frac{(1+r)^2}{r}\log\frac{N_r}{Q_r}.
\]
Differentiation gives
\[
 H'(r)=(1-r)(1+r)
 \left[-\frac1{r^2}\log\frac{N_r}{Q_r}
       +\frac{16(1+r)}{N_rQ_r}\right]\ge0.
\]
For the inequality, use $\log z\le(z-1)/\sqrt z$ and
\[
 4(1+r)^2-N_rQ_r
 =3+12r(1-r)+4r^3(6-5r)\ge0.
\]
Since $H(1)=8\log3$, this proves $g_3\le4\log3\,w$ on the left.
For general $q$, put $y=(1-r)/r\ge0$ and $k=u/\pi\ge1$. Then
\[
 \partial_q(g_q/w)
 =\frac{q[y+3+1/(y+1)]}{k^2y^2+q^2}
 \le\frac{q(y+4)}{y^2+q^2}\le\frac12+\frac4q.
\]
The last step follows from $2qy\le y^2+q^2$. Integration from $1$
to $q$ proves \eqref{eq:kernel-variance}. Finally, if $r\ge1/3$,
then $p\ge1/4$ and $g_q\le\log q$. If $r\le1/3$, then
$y=1/r-1\ge2$, and $\log(1+t^2)\le t$ gives
\[
 \frac{g_q}{p}\le\frac{y+2}{2}\log(1+q^2/y^2)
 \le\frac{q(y+2)}{2y}\le q.
\]
This proves \eqref{eq:kernel-mean}.

Put $C_q=4\log q+(q-1)/2$ and $M_q=\max(q,4\log q)$.
Equations~\eqref{eq:bulk-mode-mean}, \eqref{eq:negative-exact-ratio}
and Lemma~\ref{lem:variance} yield
\begin{align}
 \log R_3&\le-n\left(1-\frac1{Y-\log2}\right)D_3
                    +\frac{32n\log3}{Y^2},\label{eq:negative-three-bound}\\
 \log R_q&\le-n\left(1-\frac1{Y-\log2}\right)D_q
                    +\frac{8nC_q}{Y^2},\label{eq:negative-var-bound}\\
 \log R_q&\le-(n+1)D_q+M_q\left(\frac n{Y-\log2}+1\right).
                    \label{eq:negative-mean-bound}
\end{align}

For $0\le u\le\pi$, use $B=(1+e^u)/\pi$ instead. Bound
\eqref{eq:bulk-mode-mean} remains valid: at the midpoint the
consecutive probability ratio is less than one because
$B\ge2/\pi$ and $n\ge250$; log-concavity therefore places the mode
before the midpoint. The kernel becomes
\[
 \frac12\log\frac{(u/\pi-1/r)^2+q^2}{(u/\pi-1/r)^2+1},
\]
and \eqref{eq:kernel-mean} still holds. Its proof for $r\le1/3$ uses
$|u/\pi-1/r|\ge1/r-1\ge2$; its other case is unchanged.
Thus \eqref{eq:negative-mean-bound} applies in this range as well.

\subsection{Summation of the negative-side pole bounds}

On $0\le u\le\pi$, the rational interval checks specified below
show that the bounds for $R_3,R_5,R_7,R_9$ are all less than $e^{-12}$.
For $Q=9$, the exponent on the right of
\eqref{eq:universal-pole-tail}, before its factor $Q/2$, is also less
than $-12$. Hence the whole sum is less than $.001$.

On $\pi\le u\le50$, the same checks give
$R_3<.79$, $R_5<e^{-4}<.02$, and bounds $e^{-12}$ for the remaining
finite poles and the tail exponent. There are at most $23$ finite
terms with $q\ge7$, and $Q/2\le25.5$, so the total is less than $.811$.

It remains to treat $50\le u\le4\sqrt n$. Here
\[
 Y\ge u+1.37,\quad Y-\log2\ge u+.67,\quad
 1-1/(Y-\log2)>.96,\quad n/u^2\ge1/16.
\]
The first inequality follows from
$n/u\ge\sqrt{250}/4>e^{1.37}$. Using
$\log(1+x)\ge2x/(2+x)$, the bound for $q=3$ reduces to
\begin{equation}\label{eq:negative-rational-polynomial}
 \left(1-\frac1{u+.67}\right)
       \frac{39.44u^2}{u^2+49.35}
 -\frac{35.156u^2}{(u+1.37)^2}\ge3.9.
\end{equation}
After multiplication by the positive common denominator, the
numerator has positive coefficients in $u-50$. Thus
$R_3\le e^{-3.9/16}<.79$.
For $q=5$,
\[
 -\log R_5\ge\frac n{u^2}
 \left\{\frac{.96\cdot9.86\cdot24}
 {2+9.87\cdot26/2500}-8C_5\right\}>2.53,
\]
so $R_5<.08$.
Let $Q$ be the smallest odd integer at least $u$. For $7\le q\le Q$,
$u\ge\max(50,q-2)$, and the following rational-polynomial comparison
holds:
\[
 \frac{.96\cdot9.86(q^2-1)}
 {2+9.87(q^2+1)/\max(50,q-2)^2}
 -32\log q-4(q-1)\ge16q.
\]
For completeness, replace $\log q$ by its tangent bound
$1.946+(q-7)/7$. With
\[
 A=.96\cdot9.86,\quad D=9.87,\quad
 L(q)=\frac{172}{7}q+26.272,
\]
the required numerators are
\[
 2500A(q^2-1)-L(q)[5000+D(q^2+1)]\quad(7\le q\le52)
\]
and
\[
 A(q^2-1)(q-2)^2
 -L(q)[2(q-2)^2+D(q^2+1)]\quad(q\ge52).
\]
The first has positive Bernstein coefficients on $[7,52]$ and the
second has positive coefficients in $q-52$. Therefore $R_q\le e^{-q}$,
and the sum from $q=7$ through $Q$ is less than $.0011$.

At $Q$, one has $D_Q>1.18$ and
\[
 \frac{8C_Q}{Y^2}
 \le\frac{32\log(u+2)+4(u+1)}{(u+1.37)^2}<.3.
\]
The last expression decreases for $u\ge50$, and its value at $50$
is less than $.3$. Hence $R_Q<e^{-.8n}$. The extra exponent in
\eqref{eq:universal-pole-tail} is at most $.051n$, and
$Q/2<3\sqrt n$. The infinite tail is less than $.001$.
Altogether,
\[
 \sum_{q\ge3,\ q\text{ odd}}R_q
 <.79+.08+.0011+.001<.88.
\]
This establishes all the quantitative bounds in Lemma~\ref{lem:bulk}.
At a phase node, the nearest conjugate pair contributes
$2C_s(-1)^j|\mathcal R_{n,s}(s+i\pi)|$; the remaining pairs have
strictly smaller total modulus. This proves its sign conclusion.

\subsection{Exact verification of the scalar inequalities}

We specify the finite checks so that the numerical constants in the
proof are reproducible independently of root computations. Every
terminating decimal in this section denotes the corresponding rational
number. The elementary bounds $3.14159<\pi<355/113$ suffice.
For $x\ge0$ and $x<m+2$, use
\[
 E_m(x)=\sum_{k=0}^m\frac{x^k}{k!},\qquad
 E_m(x)\le e^x\le
 E_m(x)+\frac{x^{m+1}}{(m+1)!}\frac1{1-x/(m+2)}.
\]
Take $m=100$ in the variance checks, $m=80$ in the bounded-interval
pole checks, and $m=220$ for the final tail exponent at $n=250$.
For logarithms, first reduce the argument to $[1,2]$ by powers of
$2$ and use, with $z=(x-1)/(x+1)$,
\[
 2\sum_{k=0}^{39}\frac{z^{2k+1}}{2k+1}
 \le\log x\le
 2\sum_{k=0}^{39}\frac{z^{2k+1}}{2k+1}
 +\frac{2z^{81}}{81(1-z^2)}.
\]
Reciprocation handles arguments below one.
For $I_0(v)$ on $[.7,1]$, retain its first $25$ terms and bound its
remaining terms by the geometric ratio $(v^2/4)/26^2$.

The following explicit upper bounds control the two series used in
the variance proof. Put $r=e^{-8t}$ and define
\begin{align*}
 \overline S(t)&=1+4e^{-2t}+9e^{-6t}
 +e^{-12t}\left\{\frac{16}{1-r}
 +\frac{8r}{(1-r)^2}+\frac{r(1+r)}{(1-r)^3}\right\},\\
 \overline T(t)&=e^{-2t}+4e^{-6t}
 +e^{-12t}\left\{\frac9{1-r}
 +\frac{6r}{(1-r)^2}+\frac{r(1+r)}{(1-r)^3}\right\}.
\end{align*}
They follow from $j(j-1)\ge12+8(j-4)$ for $j\ge4$.
Replacing exponentials by the rational bounds above makes every entry
of Table~\ref{tab:variance-intervals} a rational comparison. In its last
row, $\overline\Delta$ is obtained from $\Delta$ by replacing $S,T$
with $\overline S,\overline T$.

\begin{table}[ht]
\centering
\caption{Finite scalar checks for Lemma~\ref{lem:variance}.}
\label{tab:variance-intervals}
\begin{tabular}{p{.30\textwidth}p{.62\textwidth}}
\toprule
Intervals or parameters & Verified inequality\\
\midrule
$[a,b]=[.7+i/1000,.7+(i+1)/1000]$, $0\le i<300$
& $b[e^{-a}I_0(a)]^2<.47^2$\\[3pt]
$[a,b]$ one of $(2,3)$, $(3,4)$, $(4,6)$, $(6,8)$, $(8,10)$, $(10,11)$
& $\{.47[\overline S(a/7)+e^{-a/2}\overline S(a/3)+.003]\}^2<(8/b)^3$\\[3pt]
$[a,b]=[i/2,(i+1)/2]$, $22\le i<34$
& $G=8/b-1/4-\overline\Delta(a)>0$ and $\overline\Delta(a)<G^2$\\
\bottomrule
\end{tabular}
\end{table}

Here are the bounded-interval pole checks in equally explicit form.
For a box $[a,b]$, use
\[
 D_q^-:=\frac12\log\left(1+
       \frac{(q^2-1)9.86}{b^2+9.87}\right),\qquad
 \ell_2=.7>\log2.
\]
On $0\le u\le\pi$, take $[a,b]=[i/20,(i+1)/20]$,
$0\le i<63$, and
$Y^- =\log[250(1+e^a)/(355/113)]$.
For $q=3,5,7,9$, set
\[
 f_q=-D_q^-+\frac{M_q}{Y^--\ell_2},\qquad
 c_q=M_q-D_q^-.
\]
The checks give $f_q<0$ and $250f_q+c_q<-12$.
For $q=9$, they also give
\[
 f_q+\tau<0,\qquad 250(f_q+\tau)+c_q+\tau<-12,
 \quad \tau=\frac{b^2}{2q^2\,9.86}.
\]

On $\pi\le u\le50$, take
$a=3.14+i/10$, $b=\min(a+.1,50)$, $0\le i<469$, and
$Y^-=a+\log(250/b)$. Bound $q=3$ by
\eqref{eq:negative-three-bound}, using $D_3^-$, $Y^-$ and $.7$;
the resulting exponent at $n=250$ is less than $-39/160$.
Let $Q$ be the least odd integer at least $\max(5,b)$.
For each $q=5,7,\ldots,Q$, use the smaller of the mean bound and the
variance bound, with the same substitutions. The exponent is less
than $-4$ for $q=5$ and less than $-12$ for $q\ge7$.
For $q=Q$, adding $(n+1)\tau$ still gives an exponent below $-12$.
In each case the coefficient of $n$ is strictly negative and decreases
as $Y$ increases. Since $Y_n=Y_{250}+\log(n/250)$, the $n=250$ check
therefore proves the assertion for every $n\ge250$.

The remaining checks are polynomial positivity in
\eqref{eq:negative-rational-polynomial} and in the two displayed
polynomials in $q$, together with their scalar endpoint comparisons.
For the cubic on $[7,52]$, substitute $q=7+45t$ and convert its power
coefficients $a_k$ to Bernstein coefficients
$\sum_{k\le i}a_k\binom ik/\binom3k$; all four are positive.
For the unbounded intervals, direct expansion in $u-50$ or $q-52$
has strictly positive coefficients. Thus the finite verification uses
rational arithmetic throughout.

\section{The small negative zeros}

Throughout this section $n\ge250$, $p=n+1$ and $K=\lfloor\sqrt n/3\rfloor$. Set
\[
c(t)=\frac{t-1}{\log t}\quad(t>1),\qquad c(1)=1,\qquad \ell(t)=\log c(t).
\]

\subsection{Coefficient estimates}
For $t\ge1$,
\begin{equation}\label{eq:logmean-curvature}
\ell'(t)\ge\frac1{2t},\qquad -\ell''(t)\ge\frac5{12t^2}.
\end{equation}
For the first inequality, put $u=\log t$ and use
$t\ell'(t)=1/(1-e^{-u})-1/u\ge1/2$.
For the second, multiplication by $12u^2(e^u-1)^2$ reduces the claim to
\[
(7u^2-12u-12)e^{2u}+2(5u^2+12u+12)e^u-(5u^2+12u+12)\ge0.
\]
Its coefficients through degree four vanish. For $h\ge5$, its coefficient multiplied by $h!$ is
\[
2^{h-2}(7h^2-31h-48)+10h^2+14h+24>0:
\]
at $h=5$ this is $120$, and for $h\ge6$ both displayed terms are positive. This proves \eqref{eq:logmean-curvature}.

Uniformly for the integers $1\le j\le n$,
\begin{equation}\label{eq:eulerian-evaluation}
A_n(j)=n!c(j)^p(1+e_{n,j}),\qquad |e_{n,j}|<2^{-120}.
\end{equation}
The assertion is exact at $j=1$. For $j\ge2$, the Frobenius identity, reciprocity and the partial fractions of $1/(1-e^{-z})$ give
\[
\frac{A_n(j)}{n!c(j)^p}
=\sum_{h\in\mathbb Z}\left(1+\frac{2\pi i h}{\log j}\right)^{-p}.
\]
Put $a=2\pi/\log n$. The nonzero terms have total modulus at most
\begin{equation}\label{eq:evaluation-tail}
2(1+a^2)^{-p/2}
\left(1+\frac1{pa^2/(1+a^2)-1}\right).
\end{equation}
Indeed the decreasing-integral bound applies after using
$1+a^2t^2\ge(1+a^2)t^{2a^2/(1+a^2)}$ for $t\ge1$.
Both $p\log(1+a^2)$ and $pa^2/(1+a^2)$ increase with $n\ge250$.
At $n=250$, $\pi>3$ and $\log250<6$ imply $a^2>1$ and
$pa^2/(1+a^2)>125$. Thus \eqref{eq:evaluation-tail} is less than
$4\cdot2^{-125}<2^{-120}$.

Summing \eqref{eq:bonferroni} against $A(n,h)$ gives
\begin{equation}\label{eq:B-small}
A_n(j)-(n+1)A_n(j-1)\le B(n,j)\le A_n(j).
\end{equation}
For $j\ge2$, \eqref{eq:logmean-curvature} gives
$(c(j-1)/c(j))^p\le e^{-p/(2j)}$; at $j=1$, $B(n,1)=n!$ exactly. Consequently
\begin{equation}\label{eq:B-small-relative}
\left|\frac{B(n,j)}{n!c(j)^p}-1\right|<\frac1{1000}
\qquad(1\le j\le K+4),
\end{equation}
and $0<B(n,j)/(n!c(j)^p)<2$ for every $1\le j\le n$.
For completeness, the error in \eqref{eq:B-small-relative} is at most
\[
2^{-120}+(n+1)(1+2^{-120})
\exp\left(-\frac{n+1}{2(\sqrt n/3+4)}\right).
\]
This decreases for $n\ge250$: after setting $t=\sqrt n$, the relevant logarithmic derivative is
\[
\frac{3(t^2+24t-1)}{2(t+12)^2}-\frac{2t}{t^2+1}>0.
\]
At $n=250$ the exponent exceeds $753/56>13$, and $251e^{-13}<1/1000$ with room for the $2^{-120}$ term.

\subsection{Alternating signs}
For $2\le k\le K+1$ define
\[
r_{n,k}=\left(\frac{c(k-1)}{c(k+1)}\right)^{p/2},
\qquad s_k=-\log r_{n,k}.
\]
Then
\begin{equation}\label{eq:small-separator-sign}
\sgn P_n(-r_{n,k})=(-1)^k.
\end{equation}
To prove this, let $\delta_k=2\ell(k)-\ell(k-1)-\ell(k+1)$.
Equation \eqref{eq:logmean-curvature} yields
$\delta_k\ge5/[12(k+1)^2]$ and $p\delta_k\ge2$.
For $K=5$, use $p\ge251$ and $k+1\le7$; for $K\ge6$, use
$n\ge9K^2$ and $(15/4)(K/(K+2))^2\ge135/64>2$.

Set $r=r_{n,k}$ and $T_j=c(j)^p r^j$. Strict log-concavity gives
\[
\frac{T_{k-h}}{T_k},\frac{T_{k+h}}{T_k}\le q^h,\qquad
q=e^{-p\delta_k/2}\le e^{-1}<\frac38.
\]
Each one-sided alternating tail is decreasing, so
$\sum_{j=1}^n(-1)^{j-k}T_j/T_k\ge1-2q>1/4$.
Replacing the ideal coefficients by $B(n,j)/n!$ changes this by less than
\[
\frac1{1000}\left(1+\frac{2q}{1-q}\right)
+\frac{q^4}{1-q}
<\frac{11}{5000}+\frac{81}{2560}<\frac14.
\]
Here \eqref{eq:B-small-relative} applies through $j=k+3$, and the global ratio in $(0,2)$ applies thereafter. This proves \eqref{eq:small-separator-sign}. The $s_k$ strictly decrease with $k$.

\subsection{Joining these signs to the bulk}
\begin{lemma}\label{lem:small}
There is a phase node $s_*$ with
\[
\phi(s_*)=-(K+1)\pi,\qquad 0<s_*<3\sqrt n,\qquad s_*<s_K.
\]
The phase-node sign followed by the signs at $s_K,s_{K-1},\ldots,s_2$ and near $x=0-$ supplies $K$ distinct negative roots after this node.
\end{lemma}
\begin{proof}
For $s>0$, write $C=1+e^{-s}$, $a_i=C/\beta_i$, $t=C/s$, and
$\mu(t)=tF_n'(t)/F_n(t)$. Directly from \eqref{eq:phase},
\begin{equation}\label{eq:positive-phase-mean}
-\phi(s)=2\arctan(\pi/s)
+\sum_i\arctan\frac{a_i\pi}{s(s+a_i)+\pi^2}
<\frac{\pi}{s}\{\mu(t)+1\}.
\end{equation}
For $0<t\le1$, $\mu(t)\le n-1$. In fact the only negative coefficient of
$(n-1)F_n-tF_n'$ is $-n!t^n$, and the term of degree $n-2$ covers it:
$A(n,3)\ge n(n-1)$ for $n\ge5$. The latter starts at $A(5,3)=66$ and follows inductively from $A(n+1,3)\ge3A(n,3)$. Thus
\[
-\phi(3\sqrt n)<\pi\sqrt n/3<(K+1)\pi.
\]

To compare with $s_K$, give $U>0$ the density proportional to
$g(u)e^{-su}$, where $g(u)=A_n(Cu)$. Termwise integration gives
\[
\E U=\frac{\mu(C/s)+1}{s}=:M.
\]
The function $h(u)=ug'(u)/g(u)$ is concave, since the negative roots of $A_n$ express it as a sum of terms $u/(u+\alpha)$, $\alpha\ge0$. Integration by parts and Jensen's inequality yield
$sM=1+\E h(U)\le1+h(M)$.
As $1/u+g'(u)/g(u)$ strictly decreases, it follows that
\begin{equation}\label{eq:mean-saddle}
s>\frac1b+\frac{g'(b)}{g(b)}\quad\Longrightarrow\quad M<b.
\end{equation}
Take $s=s_K$ and $b=K+1$. The curvature bound gives
\[
s-p\{\ell(b)-\ell(b-1)\}=\frac p2\delta_K\ge1.
\]
Concavity of $\log A_n$ and \eqref{eq:eulerian-evaluation} imply
\[
\frac{g'(b)}{g(b)}
\le C\{\log A_n(b)-\log A_n(b-1)\}
\le C\{p(\ell(b)-\ell(b-1))+4\cdot2^{-120}\}.
\]
Here $b\ge6$, $b\le(2/5)\sqrt n$, $s\ge p/(2b)\ge(5/4)\sqrt n$, and
$p(\ell(b)-\ell(b-1))\le p$. Hence the extra factor $C-1=e^{-s}$ changes the last bound by less than $1/1000$; the function
$(n+1)e^{-(5/4)\sqrt n}$ decreases and is already below $1/1000$ at $250$.
Since $1/b\le1/6$, \eqref{eq:mean-saddle} gives $M<b$. Equation
\eqref{eq:positive-phase-mean} now gives $\phi(s_K)>-(K+1)\pi$.

Finally put $t_0=2/\pi$. Lemma~\ref{lem:weightedmean} and
$\arctan x\ge x/(1+x)$ give
\[
-\phi(0)=\pi+\sum_i\arctan(t_0/\beta_i)
\ge\pi+\mu(t_0)-1
\ge\pi+(n-1)/2>(K+1)\pi.
\]
Strict phase monotonicity proves the three assertions about $s_*$. Its sign is $(-1)^{K+1}$ by Lemma~\ref{lem:bulk}; the subsequent separator signs alternate, and the sign near zero is negative because $B(n,1)=n!$. Their $K$ intervening intervals are disjoint.
\end{proof}

\section{Pairs of large negative zeros}

For an integer $\ell\ge1$ write
\[
\rho_\ell=\frac{\ell+1}{\ell},\quad
a_\ell=\ell^{1-\ell},\quad
b_\ell=(\ell+1)^{1-\ell}+\ell^{-\ell},\quad
c_\ell=(\ell+1)^{-\ell},
\]
and set
\begin{equation}\label{eq:outer-separators}
\lambda_\ell^-=\frac{c_\ell}{2b_\ell},\qquad
\lambda_\ell^0=\frac{b_\ell}{2a_\ell},\qquad
\lambda_\ell^+=\frac{2b_\ell}{a_\ell},\qquad
X_\ell^\nu=\lambda_\ell^\nu\rho_\ell^n.
\end{equation}

\begin{lemma}\label{lem:outer}
Let $n\ge250$ and $L=\lfloor\sqrt{n/(4\log n)}\rfloor$. For $1\le\ell\le L$,
\[
\sgn P_n(-X_\ell^-)=\sgn P_n(-X_\ell^+)=(-1)^n,\qquad
\sgn P_n(-X_\ell^0)=(-1)^{n+1}.
\]
The intervals $(X_\ell^-,X_\ell^+)$ are disjoint. Consequently $P_n$ has at least $2L$ distinct negative zeros with moduli greater than $X_L^-$.
\end{lemma}
\begin{proof}
Fix $\ell$ and set $X=\lambda\rho_\ell^n$,
$Z=(\ell+1)^{2n}X^{-2\ell}$. The four ideal weights
$W_u(k,h)$ with $(k,h)\in\{\ell,\ell+1\}^2$ have signed sum and absolute sum
\[
ZQ(\lambda),\quad ZM(\lambda),\qquad
Q=a_\ell\lambda^2-b_\ell\lambda+c_\ell,\quad
M=a_\ell\lambda^2+b_\ell\lambda+c_\ell.
\]
Throughout $\lambda_\ell^-\le\lambda\le\lambda_\ell^+$,
\begin{equation}\label{eq:lambda-bounds}
\frac1{2(\ell+4)}<\lambda\le4.
\end{equation}
Indeed $b_\ell/c_\ell=\ell+1+\rho_\ell^\ell<\ell+4$, while
$b_\ell/a_\ell\le2$.

For fixed $h$, consecutive ideal $k$ ratios are
$W(k+1,h)/W(k,h)=\rho_k^n/(hX)$.
The outward right and left ratios beyond the two central columns are therefore at most
\[
q_R=2(\ell+4)e^{-n/(\ell+1)^2},
\qquad q_L=4n e^{-n/\ell^2},
\]
respectively; a left tail is absent at $\ell=1$.
The left estimate uses $h\le n$, sufficient for the actual triangular sum.
In either central column, the corresponding outward $h$ ratios are at most
$q_R$ and $16e^{-n/\ell^2}$.

Put $H=\sqrt{n/(4\log n)}\ge3$. Then $\ell+1\le4H/3$, $\ell+1\le\sqrt n/3$, and
\[
\frac n{(\ell+1)^2}\ge\frac94\log n,\qquad
\frac n{\ell^2}\ge4\log n,\qquad
q_R\le\frac76n^{-7/4}<\frac1{1000}.
\]
The other three ratios are also below $1/1000$. Summing the two geometric tails in each coordinate shows that the ideal mass outside the central four cells is at most
\[
\{(1+2/999)^2-1\}ZM<ZM/200.
\]
By \eqref{eq:bonferroni}, replacing the central ideal weights by the actual Eulerian products changes their sum by less than $ZM/1000$: each of the two relative errors is bounded by
$(n+1)e^{1-3\sqrt n}$, and $2(n+1)e^{1-3\sqrt n}<1/1000$.
Thus, after division by $(-X)^n Z$, the actual polynomial differs from $Q(\lambda)$ by less than $M(\lambda)/100$.

The three quadratic signs have uniform margins. Put
$u=(\ell+1)^{1-\ell}$, $v=\ell^{-\ell}$, $w=u/v$. Then
\[
\frac{uv}{a_\ell c_\ell}=\frac{\ell+1}{\ell},
\qquad w=(\ell+1)(\ell/(\ell+1))^\ell>\frac{\ell+1}{3}.
\]
For $\ell\ge2$, $w>1$, and
\[
\frac{b_\ell^2}{a_\ell c_\ell}
>\left(1+\frac1\ell\right)
 \left(\frac{\ell+1}{3}+\frac3{\ell+1}+2\right)
=\frac{(\ell+4)^2}{3\ell}\ge\frac{16}{3}.
\]
For $\ell=1$ the ratio is $8$. Substitution at the three values of $\lambda$ gives
\[
\frac{Q(\lambda^-)}{M(\lambda^-)}>\frac{32}{99},\qquad
-\frac{Q(\lambda^0)}{M(\lambda^0)}>\frac1{15},\qquad
\frac{Q(\lambda^+)}{M(\lambda^+)}>\frac13.
\]
For the middle bound, put $D=b_\ell^2/(4a_\ell)>4c_\ell/3$ and use
$(D-c_\ell)/(c_\ell+3D)>1/15$.
These margins prove the signs of the actual polynomial.
Finally
\[
\frac{X_{\ell+1}^+}{X_\ell^-}
<8(\ell+4)e^{-n/(\ell+1)^2}=4q_R<1.
\]
The two sign changes in each disjoint interval give the asserted roots.
\end{proof}

For the short intermediate range we can increase the cutoff.
\begin{lemma}\label{lem:finitebridge}
Let $250\le n\le1600$ and $L=\lceil\sqrt n/4\rceil$. The signs and disjoint intervals of Lemma~\ref{lem:outer} hold for every $\ell\le L$. Moreover, with $u=\log X_L^-$,
\begin{equation}\label{eq:finite-phase-match}
40<u<4\sqrt n,\qquad
(2L-n-1)\pi<\phi(-u)<(2L-n)\pi.
\end{equation}
\end{lemma}
\begin{proof}
Here $\ell\le10$ and
$n\ge n_\ell:=\max(250,16(\ell-1)^2+1)$.
The four outward ratios used above have the uniform bounds
\[
\begin{array}{c|cc}
 &\text{left}&\text{right}\\ \hline
k&.036&.015\\
h&.0006&.015
\end{array}
\]
and the central relative error is below $10^{-6}$.
These follow by evaluating the decreasing functions
$4n e^{-n/\ell^2}$, $2(\ell+4)e^{-n/(\ell+1)^2}$,
$16e^{-n/\ell^2}$ and $2(n+1)e^{-(n-\ell)/(\ell+1)}$
at the ten integer endpoints $n_\ell$. Exact rational bounds for these evaluations are part of Appendix~\ref{sec:verification}.
Each tail attaches to its own central row or column, so their common factors are
\[
\eta_k=\frac{.036}{1-.036},\qquad
\eta_h=\frac{.015}{1-.015},\qquad
(1+\eta_k)(1+\eta_h)-1<.054.
\]
Thus the quadratic error is less than $(.054+10^{-6})M<M/15$.
The same margins prove the three signs; disjointness follows from
$8(\ell+4)e^{-n/(\ell+1)^2}<.06$.
Since $\lambda_L^-<1$, $u<n/L\le4\sqrt n$.

We give the phase comparison in detail. For any $\ell\le L$, put
\[
D=2(\ell+1+\rho_\ell^\ell),\quad
u=n\log\rho_\ell-\log D,\quad
z=-u+i\pi,\quad \theta=\pi-\delta,\quad
\delta=\arctan(\pi/u),\quad y=\frac z{1+e^u}.
\]
The same ten endpoint checks give
\begin{equation}\label{eq:finite-phase-constants}
\begin{gathered}
u>40,\quad u>\frac{11n}{12(\ell+1)},\quad
u>\frac{n-\ell+3}{\ell+1},\\
\frac n{2(\ell+1)}>21,\qquad
\frac{(\ell+1)^2}{n}<\frac16,\\
(n+1)e^{-n/(\ell+2)}<10^{-6},\qquad e^{-u}<10^{-6}.
\end{gathered}
\end{equation}
For the first three inequalities one uses the lower bound
$u>n\log\rho_\ell-\log(2(\ell+4))$; all three differences increase with $n$.

Write $E(y)=\sum_{r=0}^{n-1}b_r y^r$,
$b_r=(n-r)!A(n,r+1)$. The first term to the left of $b_\ell y^\ell$ has ratio less than $p_0=11/20$. More generally, for every $1\le r\le\ell$,
\[
\frac{b_{r-1}}{b_r|y|}
\le\frac n{Du}\,\frac{1+e^{-u}}{1-10^{-6}}
       \left(\frac{\rho_\ell}{\rho_r}\right)^n
<p_0\left(\frac{\rho_\ell}{\rho_r}\right)^n.
\]
This uses $n-r+1\le n$ directly. Thus all further left ratios are less than
$p_0 e^{-n/\ell^2}<q_0=1/100$.
The first ratio to the right is at most
\[
\frac{2(\ell+4)(n/\ell+4)}{(n-\ell)(1-10^{-6})}
 e^{-n/(\ell+1)^2}<q_0.
\]
The rational prefactor and the exponential decrease with $n$; the endpoint check proves the bound. Since $E$ is the reversal of $F_n$, its coefficients are log-concave, so all later right ratios are no larger.
Consequently, for $J(y)=E(y)/(b_\ell y^\ell)$ on the entire semicircle of radius $|y|$,
\[
|J-1|<\frac{p_0+q_0}{1-q_0}<.57,\qquad
\sum_{r\ne\ell}|r-\ell|\frac{|b_r y^r|}{|b_\ell y^\ell|}
<\frac{p_0+q_0}{(1-q_0)^2}<.58.
\]
At angle $\theta=\pi-\delta$, $|\Im J|<.58\delta$, $\Re J>.43$, and hence
$|\arg J|<2\delta$. Control on the whole semicircle fixes the continuous argument without a multiple of $2\pi$. Reversal now gives
\[
\phi(-u)=-(n-\ell+1)\theta+\arg J.
\]
The upper phase bound follows from
$(n-\ell+3)\delta/\pi<(n-\ell+3)/u<\ell+1$.
For the lower bound, the trapezoid estimate for $1/t$ and $D>4$ yield
\[
n-\ell-1-\ell u
\ge\frac n{2(\ell+1)}-\ell-1+\ell\log D
>\frac n{2(\ell+1)}-1>20.
\]
On the other hand,
$(n-\ell-1)\pi^2/(3u^2)<48\pi^2(\ell+1)^2/(121n)<1$.
Using $\arctan t>t-t^3/3$ gives
$(n-\ell-1)\delta/\pi>\ell$. The two bounds imply
$(2\ell-n-1)\pi<\phi(-u)<(2\ell-n)\pi$.
Taking $\ell=L$ proves \eqref{eq:finite-phase-match}.
\end{proof}

\section{The transition on the large-root side}\label{sec:bridge}

Put $X=e^u>1$ and
\[
 W_u(k,h)=k^n h^{n-k+1}X^{-(k+h-2)},\qquad k,h\ge1.
\]
The positive double sum converges, since
$W_u(k,h)\le k^nh^nX^{-(k+h-2)}$. Eulerian symmetry gives
\begin{equation}\label{eq:bridge-exact}
 \frac{P_n(-X)}{(-X)^n}
 =\sum_{\substack{k,h\ge1\\k+h\le n+1}}
 (-1)^{k+h-2}A(n,k)A(n-k+1,h)X^{-(k+h-2)}.
\end{equation}
Lemma~\ref{lem:bonferroni} bounds every absolute summand by $W_u(k,h)$.
Throughout this section, $\rho_j=(j+1)/j$.

\subsection{The path of maximizing cells}

For an integer $m\ge2$, define
\[
 f(k,h)=n\log k+(n-k+1)\log h,\qquad
 M_m=\max_{k+h=m}e^{f(k,h)}.
\]
Choose the smallest maximizing $k=k_m$, and set $h_m=m-k_m$.
For $2\le m<m+1\le\lfloor\sqrt n\rfloor$, this path increases exactly
one coordinate by one at each step. Indeed, for $h=m-k\ge2$ its
adjacent difference is
\[
 q_m(k)=n\log\rho_k-\log h-(n-k)\log\rho_{h-1}.
\]
The comparisons
\begin{align*}
 q_{m+1}(k)-q_m(k)
   &=(n-k)\log\frac{h^2}{h^2-1}-\log\rho_h>0,\\
 q_{m+1}(k+1)-q_m(k)
   &=-n\log\frac{(k+1)^2}{k(k+2)}+\log\rho_{h-1}<0
\end{align*}
follow from $n-k>h$ and $n/(k+1)^2\ge1>\log2$, respectively.
They also imply that $q_m(k)$ strictly decreases in $k$. Its zero
crossing therefore moves by zero or one position when $m$ increases.
The smallest-maximizer convention includes a tie of adjacent cells;
boundary cases use the same one-sided comparisons.

The sequence $\log M_m$ is strictly concave. Two successive steps in
one coordinate have decreasing increments by coordinate concavity.
For a $k$ step followed by an $h$ step, optimality gives
$f(K+1,H)\ge f(K,H+1)$, whereas
\[
 [f(K,H+1)-f(K,H)]
 -[f(K+1,H+1)-f(K+1,H)]=\log\rho_H>0.
\]
Thus the second increment is strictly smaller; the opposite order is
identical. Consequently the global maximizing total degree changes
one at a time as $u$ decreases, at the ordered times
$\log(M_{m+1}/M_m)$.

The two coordinate events starting from $(K,H)$ have balance times
\begin{equation}\label{eq:bridge-balances}
 \alpha=n\log\rho_K-\log H,\qquad
 \beta=(n-K+1)\log\rho_H.
\end{equation}
Consecutive events in the same coordinate are separated by more than
\begin{equation}\label{eq:bridge-spacing}
 \frac{n}{(K+1)^2},\qquad \frac{n-K+1}{(H+1)^2},
\end{equation}
respectively. Their exact differences are
$n\log((K+1)^2/(K(K+2)))+\log(H'/H)$, or
$(n-K+1)\log((H+1)^2/(H(H+2)))+(K'-K)\log\rho_{H+1}$,
where $H'\ge H$ and $K'\ge K$ along the path.

\subsection{Global descent estimates}

Fix $(K,H)$ and put $N=n-K+1$, $S=K+H$. For integer shifts
$p\ge1-K$, $q\ge1-H$, write
\[
 w_{pq}=\frac{W_u(K+p,H+q)}{W_u(K,H)},\quad
 A=\frac{\rho_K^n}{HX},\quad B=\frac{\rho_H^N}{X},\quad
 d(p)=\max\{-p,p-1,0\}.
\]

\begin{lemma}\label{lem:bridge-joint}
Suppose $K,H\ge7$, $n-S>S^2$,
\[
 \delta_k=n\log(\rho_K/\rho_{K+1})\ge12,\qquad
 \delta_h=N\log(\rho_H/\rho_{H+1})\ge12,\qquad
 e^{-2.02}\le A,B\le e^{2.02}.
\]
With $\varepsilon=10^{-4}$, one has
\[
 w_{pq}\le57AB\varepsilon^{d(p)+d(q)},\qquad
 \sum_{(p,q)\notin\{0,1\}^2}|pq|w_{pq}<.04AB.
\]
\end{lemma}

\begin{proof}
At an outside cell decrease the coordinate with the larger distance
from $\{0,1\}$. The four outward ratios are bounded by
$Ae^{-G_+(d)}$, $A^{-1}e^{-G_-(d)}$, $Be^{-J_+(d)}$, or
$B^{-1}e^{-J_-(d)}$, where
\begin{align*}
G_+(d)&=n\log\frac{\rho_K}{\rho_{K+d}}
             -\log\frac{H}{\max(1,H-d)},&&p=d+1,\\
G_-(d)&=n\log\frac{\rho_{K-d}}{\rho_K}
             -\log\frac{H+d+1}{H},&&p=-d,\\
J_+(d)&=N\log\rho_H-[N+\min(d,K-1)]\log\rho_{H+d},&&q=d+1,\\
J_-(d)&=(N-d-1)\log\rho_{H-d}-N\log\rho_H,&&q=-d.
\end{align*}
These expressions are used for admissible cells and $d\ge1$.
For example the first case has $h\ge\max(1,H-d)$; the last has
$k\le K+d+1\le S<n$.
All four losses increase with $d$. Before truncation, their successive
increments are bounded below, respectively, by
\[
 \frac{n}{(K+d+1)^2}-\log2,\quad
 \frac{n}{(K-d)^2}-\log2,\quad
 \frac{N}{(H+d+1)^2}-\log2,\quad
 \frac{n-S}{S^2}-\log2,
\]
which are positive. For the fourth, putting $m=H-d\ge2$ gives the
exact increment
$(N-d-2)\log(\rho_{m-1}/\rho_m)-\log\rho_m$.
After truncation in $G_+$ or $J_+$ monotonicity is immediate.
The first-layer losses are at least
\[
 \delta_k-\log\frac H{H-1},\quad
 \delta_k-\log\frac{H+2}H,\quad
 \delta_h-\log\rho_{H+1},\quad
 \delta_h-2\log\rho_{H-1}.
\]
Each correction is at most $1/3$. Thus every selected outward ratio
is at most $e^{2.02-12+1/3}<\varepsilon$.
This holds on the entire positive quadrant: a selected negative
$h$ step forces $k\le S<n$, whereas in a positive $h$ step increasing
$k$ only decreases the ratio, even when $n-k+1$ is negative.

Descent takes exactly $d(p)+d(q)$ steps to the central four cells.
Their largest normalized weight is
$M=\max\{1,A,B,ABH/(H+1)\}<57AB$, since $e^{4.04}<57$.
Finally
\[
 \sum_{p\in\mathbb Z}|p|\varepsilon^{d(p)}
 =\frac{1+\varepsilon}{(1-\varepsilon)^2}.
\]
Only $(1,1)$ contributes inside the central box, so the weighted tail is
less than
\[
 57AB\left\{\left[\frac{1+\varepsilon}{(1-\varepsilon)^2}\right]^2-1\right\}
 <57AB\,7\varepsilon<.04AB.
\]
\end{proof}

\begin{lemma}\label{lem:bridge-vertex}
Suppose $K,H\ge7$, $n-K-H>(K+H)^2$, and the four coordinate-neighbor
weights are at most $e^{-2}W_u(K,H)$. Then
\[
 \sum_{(k,h)\ne(K,H)}W_u(k,h)<\frac{400}{441}W_u(K,H).
\]
If also $n>1600$ and $\max(K,H)\le\sqrt n/4+2$, the sign of
$P_n(-X)$ is $(-1)^{n+K+H-2}$ and exactly $K+H-2$ zeros lie outside
$|x|=X$, counted with multiplicity.
\end{lemma}

\begin{proof}
Use the preceding descent with distances $|p|,|q|$ from one cell.
The first-layer boosts over the corresponding coordinate-neighbor
ratios are at most $H/(H-1)$, $(H+1)/H$, $\rho_H$, $\rho_{H-1}$.
Their logarithms are at most $1/6$. Subsequent losses increase by the
same four difference calculations, now starting at the single cell.
Thus $w_{pq}\le r^{|p|+|q|}$ for $r=e^{-11/6}<4/25$. Summation gives
$\sum_{(p,q)\ne(0,0)}w_{pq}<4(4/25)/(1-4/25)^2=400/441$.
The two central Bonferroni errors sum to at most
$2(n+1)e^{-(8/5)\sqrt n}<.001$ under the stated coordinate bounds.
Every other actual weight is bounded by its ideal weight. Hence the
actual central monomial of degree $n-K-H+2$ dominates the remainder
on $|x|=X$. Rouch\'e's theorem gives the count, and evaluation at $-X$
gives the sign.
\end{proof}

\subsection{The sign inside a paired transition}

Under the hypotheses of Lemma~\ref{lem:bridge-joint}, there is an exact factorization
\[
 w_{pq}=a_pb_qc_{pq},\quad
 a_p=(1+p/K)^n(HX)^{-p},\quad
 b_q=(1+q/H)^NX^{-q},\quad c_{pq}=(1+q/H)^{-p}.
\]
Let $F(u)=\sum_p(-1)^pa_p$ and $G(u)=\sum_q(-1)^qb_q$.
All sums converge absolutely. The residual is one on both axes and
$c_{11}=H/(H+1)$. Moreover
\begin{equation}\label{eq:bridge-residual}
 |c_{pq}-1|\le\frac{|pq|}{H}\max(1,c_{pq}).
\end{equation}
Indeed, if $t=1+q/H>0$ and $m=|p|$, its left side divided by the
maximum is $1-\min(t,t^{-1})^m\le m[1-\min(t,t^{-1})]\le m|t-1|$.
Consequently
\begin{equation}\label{eq:bridge-product}
 \sum_{p,q}(-1)^{p+q}w_{pq}
 =F(u)G(u)-\frac{AB}{H+1}+E,\qquad |E|<.08AB/H.
\end{equation}
Here Lemma~\ref{lem:bridge-joint} bounds the weighted tail of $w$ by
$.04AB$. The independent product $a_pb_q$ has the same bound: its
first outward ratios are at most $e^{2.02-12}<10^{-4}$ and subsequent
ratios decrease. Sum \eqref{eq:bridge-residual} outside the central box.

We also need the actual Eulerian coefficients. Suppose $n\ge1600$
and $\max(K,H)\le\sqrt n/4+2$. In the region $d(p)<H$, $d(q)<H$,
both coordinates are at most $\sqrt n/2+4$, all cells belong to
$k+h\le n+1$, and
\[
 n/k\ge(8/5)\sqrt n,\qquad (n-k+1)/h\ge(8/5)\sqrt n.
\]
The second follows at $t=\sqrt n\ge40$ from
$t^2/5-69t/10-3>0$; the first is weaker.
Bonferroni therefore bounds each near relative error by
$2(n+1)e^{-(8/5)\sqrt n}$.
Using the global envelope and
$\sum_p\varepsilon^{d(p)}=2/(1-\varepsilon)$, the normalized error
between the actual sum \eqref{eq:bridge-exact} and the infinite ideal
sum is at most
\begin{equation}\label{eq:bridge-actual-error}
 \frac{456AB}{(1-\varepsilon)^2}
 \left((n+1)e^{-(8/5)\sqrt n}+\varepsilon^H\right)
 <\frac{AB}{1000H}.
\end{equation}
For the far part use the union of $d(p)\ge H$ and $d(q)\ge H$.
After multiplying by $H/AB$, the near bound decreases for
$\sqrt n\ge40$ and is below $10^7e^{-64}$ at the endpoint; the far
bound decreases for integer $H\ge7$. This proves the strict inequality.

For two successive mode events in different coordinates of one
rectangle, the exact event times are
$\max(\alpha,\beta)$ and $\min(\alpha,\beta)-\log\rho_H$.
If their gap is less than $4$, then
$|\alpha-\beta|+\log\rho_H<4$.
There is a unique zero $\widetilde\alpha$ of $F$ within $.001$ of
$\alpha$, and a unique zero $\widetilde\beta$ of $G$ within $.001$ of
$\beta$; both functions are increasing in the respective
$2.02$ neighborhoods. To verify this explicitly, put $\eta=e^{-12}$.
Coordinate concavity gives, with missing negative indices set to zero,
\[
 a_{2+j}\le A^2\eta(A\eta)^j,\qquad
 a_{-1-j}\le A^{-1}\eta(A^{-1}\eta)^j,\qquad j\ge0.
\]
If $|u-\alpha|\le T$ and $r_T=e^T\eta$, the two tail bounds are
\begin{align*}
 \sum_{p\notin\{0,1\}}a_p
 &\le\frac{\eta(e^T+e^{2T})}{1-r_T},\\
 \sum_{p\notin\{0,1\}}|p|a_p
 &\le\frac{\eta[e^T+e^{2T}(2-r_T)]}{(1-r_T)^2}.
\end{align*}
For $T=.001$, the first is below $.00002$, whereas $1-A$ has
opposite signs of magnitude exceeding $.000999$ at the endpoints.
For $T=2.02$, the second is below $.001$, so
$F'(u)>e^{-2.02}-.001>.128$. This proves the assertions directly.
The proof for $G$ is identical. These are exact Eulerian zeros:
the Frobenius identity expresses $F$ as the nonzero prefactor
$(-1)^K(HX)^K/K^n$ times
$A_n(-1/(HX))/(1+1/(HX))^{n+1}$.

Choose $u_*=(\widetilde\alpha+\widetilde\beta)/2$.
Its distances from the two balances are below $2.001$, and
$F(u_*)G(u_*)\le0$. Equations~\eqref{eq:bridge-product} and
\eqref{eq:bridge-actual-error} show that the normalized actual sum is
less than
\[
 -\frac{AB}{H}\left(\frac H{H+1}-.08-.001\right)<-.79AB/H<0.
\]
This sign is opposite to the two separators around the paired cluster.
The midpoint lies between those separators: the independent zeros
are within $.001$ of their balances, and the separators are at least
$2$ from the actual events.

\subsection{The phase at a separator}

\begin{lemma}\label{lem:bridge-phase}
Let $n>1600$, $u\ge3\sqrt n$, and let $(K,H)$ maximize $W_u$, with
$K,H\ge7$ and coordinate-neighbor ratios at most $e^{-2}$.
For $r=K+H-2$,
\[
 |\phi(-u)-(r-n)\pi|<3\pi/4.
\]
\end{lemma}

\begin{proof}
Set $C=1+e^u$, $z=-u+i\pi$, $N=n-K+1$, and
$H_n(y)=\sum_{\ell=0}^{n-1}A(n,\ell+1)y^\ell/(n)_\ell$.
Eulerian symmetry gives $R_{-u}(z)=n!C^nz^{-n-1}H_n(z/C)$.
By Lemma~\ref{lem:factorial}, $H_n$ has negative real roots, so its
coefficients tilted by a positive geometric sequence are log-concave.
Put $\varrho=|z|/C$ and
$c_\ell=A(n,\ell+1)\varrho^\ell/(n)_\ell$.

Comparison first in $k$ and then in $h$ gives $K,H<n/u+1$ and
\begin{equation}\label{eq:bridge-h-mode}
 N\log\rho_H\le u\le N\log\rho_{H-1}.
\end{equation}
The Eulerian relative error at $K$ is below $.001$, since
$n/K>3\sqrt n-9$ and $(n+1)e^{9-3\sqrt n}<.001$.
Writing $A_\pm$ for the $k$-neighbor ratios, we have
\begin{align*}
 \frac{c_K}{c_{K-1}}
 &\le\frac{A_+}{.999}\frac{H\sqrt{u^2+\pi^2}}N\frac XC,\\
 \frac{c_{K-2}}{c_{K-1}}
 &\le\frac{A_-}{.999}\frac{(N+1)C}{HX\sqrt{u^2+\pi^2}}.
\end{align*}
Here $u>120$, $N>1585$, and \eqref{eq:bridge-h-mode} gives
$Hu/N\le7/6$ and $(N+1)/(Hu)<1.001(8/7)$.
Also $C/X<1.001$ and $\sqrt{1+\pi^2/u^2}<1.001$.
Both ratios are therefore below $1.18e^{-2}<1/6$.
Log-concavity bounds the sum of the other terms by
$2\sum_{j\ge1}6^{-j}=2/5$ times the term of degree $K-1$.

That monomial dominates on $|y|=\varrho$. Along the upper semicircle
the remaining factor stays within distance $.4$ of $1$, so its
canonical argument has magnitude below $.42$. It starts positive,
which fixes the argument without a winding ambiguity. With
$\theta=\arctan(\pi/u)$ we obtain
\[
 \phi(-u)=-(N+1)(\pi-\theta)+\zeta,\qquad |\zeta|<.42.
\]
The logarithmic-mean bounds
$1/\log\rho_{H-1}\ge\sqrt{H(H-1)}>H-.6$ and
$1/\log\rho_H\le H+.5$ give $H-.6<N/u\le H+.5$.
Furthermore
\[
 0<(N+1)/u-(N+1)\theta/\pi
 <(n+1)\pi^2/(3u^3)<.004,\qquad 1/u<.009.
\]
Thus $H-.604<(N+1)\theta/\pi<H+.509$.
Subtracting $(r-n)\pi=-(N+1)\pi+H\pi$ bounds the phase error by
$.604\pi+.42<3\pi/4$.
\end{proof}

\subsection{Assembling the transition}

\begin{lemma}\label{lem:bridge}
For every $n>1600$ there are an integer $r<3\sqrt n/5$ and a number
$u\in[4\sqrt n-8,4\sqrt n]$ such that exactly $r$ zeros of $P_n$ lie
outside $|x|=e^u$, all are distinct and negative, and
\[
 \operatorname{sgn}P_n(-e^u)=(-1)^{n+r},\qquad
 |\phi(-u)-(r-n)\pi|<3\pi/4.
\]
\end{lemma}

\begin{proof}
Write $v=\sqrt n>40$ and $L=\lfloor\sqrt{n/(4\log n)}\rfloor$.
Lemma~\ref{lem:outer} supplies $2L$ distinct negative roots beyond
its innermost separator $X_L^-$.
Start at $u_{\rm out}=\log X_L^- -2$, with central cell $(L+1,L+1)$.
Its left ratios are below $e^{-2}/2$ and its right ratios below
$e^2/1000<e^{-2}$, by the outer lemma.
Since $\log n>25/4$, we have $7\le L\le v/5$ and $L+1\le9v/40$.
For $D=2(L+1+\rho_L^L)<2(L+4)$,
\[
 u_{\rm out}=n\log\rho_L-\log D-2
 >40v/9-\log(2v/5+8)-2>37v/9>4v.
\]
The logarithm comparison holds at $v=40$ and its right side has the
larger derivative thereafter. Lemma~\ref{lem:bridge-vertex} gives an
exact exterior count of $2L$, already supplied by the outer lemma.

We verify all geometry down to $u=4v-8$. A maximum of $W_u$ exists
on the infinite quadrant. Comparison first in $k$ gives $K<n/u+1<n$,
and then comparison in $h$ gives $H<n/u+1$.
At an event its stepping coordinate is less than $n/u$, directly
from \eqref{eq:bridge-balances}. Put $a=v^2/(4v-8)$.
Thus the stepping coordinate is less than $a$ and the other less
than $a+1$. The same-coordinate spacings exceed
\[
 \frac{v^2}{(a+1)^2},\qquad \frac{v^2-a}{(a+1)^2},
\]
both greater than $12$. For $v\ge41$, clearing denominators in the
second comparison gives
$4\{v^3(v-41)+18v^2+192v-192\}>0$.
For $40<v<41$, $a<11$, so the stepping coordinate is at most $10$
and the other at most $11$; the smaller spacing exceeds $1590/121>13$.
In a paired rectangle both stepping coordinates are less than $a$,
so its $\delta_k,\delta_h$ exceed $12$ as well.
Every mode center satisfies
\[
 \max(K,H)<a+1<v/4+2,\quad
 K+H<v/2+4<3v/5,\quad n-K-H>(K+H)^2.
\]
The mode path applies since its total degree remains below
$\lfloor\sqrt n\rfloor$. Its coordinates never decrease as $u$
decreases, and hence remain at least $L+1\ge8$.

Join consecutive events into a cluster when their gap is \emph{less}
than $4$; gaps at least $4$ separate clusters.
A cluster contains at most two events: three would put two
same-coordinate events within $8$. A pair uses both coordinates.
The spacing lower bounds depend only on the earlier event, so this
also excludes a hypothetical third event just below the eight-unit band.

Choose $u_{\rm in}\in[4v-8,4v]$ at distance at least $2$ from every
event. If $4v$ is already safe, use it. Otherwise take the cluster
with an event within $2$ of $4v$ and move to distance $2$ below its
lowest event. Its width is less than $4$, so the displacement is
less than $8$; the next cluster is at least $4$ away.
Use the two endpoints and the midpoints of intervening separating
gaps as separators. A coordinate neighbor of their maximizing cell
belongs to an adjacent total degree, so its ratio is at most $e^{-2}$.
Lemma~\ref{lem:bridge-vertex} applies at each separator.

Across a singleton cluster the total degree increases by one.
The opposite separator signs give one negative real zero.
Across a paired cluster it increases by two; the separator signs
agree and the intermediate sign established above gives two distinct
negative real zeros. The Rouch\'e counts give exactly one or two
zeros in the respective annuli. Starting with the $2L$ outer zeros,
all $r=K+H-2$ exterior zeros at the final separator are therefore
distinct and negative.
Finally $u_{\rm in}\ge4v-8>3v$, so Lemma~\ref{lem:bridge-phase}
gives the phase estimate, while the coordinate bound gives
$r<3v/5$.
\end{proof}

All decimal constants here are terminating rational numbers.
Their exponential comparisons and polynomial endpoint checks are also
certified by the accompanying standard-library verification scripts.


\section{Completion of the root count}\label{sec:count}
We now prove Theorem~\ref{thm:main}. The coefficients satisfy
\[
B(n,0)=0,\qquad B(n,1)=n!,\qquad B(n,n)=1,
\qquad B(n,k)>0\quad(1\le k\le n).
\]
Thus zero is simple and the degree is $n$. It suffices to construct $n-1$ distinct negative zeros. The exact certificates in Appendix~\ref{sec:verification} do this for $n\le250$.

Suppose first that $250\le n\le1600$, put $L=\lceil\sqrt n/4\rceil$, and let $u=\log X_L^-$ be the endpoint in Lemma~\ref{lem:finitebridge}. There are $2L$ distinct negative zeros with modulus greater than $e^u$, and $P_n(-e^u)$ has sign $(-1)^n$. Moreover
\[
0<u<4\sqrt n,\qquad
\phi(-u)<(2L-n)\pi<(2L+1-n)\pi.
\]
Put $r=2L$ and $j=r+1-n$. Since the phase is strictly increasing, its node of index $j$ occurs after $s=-u$. Its sign is opposite to that of the endpoint. It therefore supplies one additional zero.

For $n>1600$, Lemma~\ref{lem:bridge} supplies a separator $u\in[4\sqrt n-8,4\sqrt n]$ and an integer $r$ such that exactly $r$ negative zeros lie outside its circle, the endpoint sign is $(-1)^{r-n}$, and
\[
\phi(-u)<(r+1-n)\pi.
\]
Again the node $j=r+1-n$ occurs after the separator and supplies one additional zero. In both cases this node is in the bulk range of Lemma~\ref{lem:bulk}.

In either range let $K=\lfloor\sqrt n/3\rfloor$. The bounds on $r$ in the two lemmas give $r+K+2<n$, so $j<-(K+1)$. Continue through the consecutive phase nodes
\[
j,j+1,\ldots,-(K+1).
\]
Their alternating signs yield $n-r-K-2$ more zeros. Lemma~\ref{lem:small} provides $K$ further zeros after the last node. All intervals are disjoint, and their number is
\[
r+1+(n-r-K-2)+K=n-1.
\]
Together with the simple zero at zero, these exhaust the degree. Each is simple, because $n$ distinct zeros have been obtained. This proves the theorem.

\appendix
\section{Exact certificates and reproducibility}\label{sec:verification}

The finite row certificates and the scalar inequalities in the analytic
proof are checked using integers and rational numbers. The supplied
verifier uses Python's standard library; it requires no computer algebra
system, numerical root finder, or external package. From the supplement's
root directory, the command is
\begin{verbatim}
python src/verify_all.py
\end{verbatim}
The script also accepts invocation by its absolute path from another
working directory. All component paths are resolved relative to the
script. It checks that the row certificate ends at $250$, checks the
presence of every required verifier, and stops at the first failed
assertion or process. Python must be run without the optimization option
\texttt{-O}, which disables assertions.

\subsection{The finite row certificates}

For $1\le n\le250$, the file
\path{certificates/sign_certificates_n1_250.json} contains $n-1$ rational
numbers
\[
 0>q_{n,1}>q_{n,2}>\cdots>q_{n,n-1}.
\]
The verifier independently computes the Eulerian numbers from
\[
 A(n,k)=\sum_{i=0}^k(-1)^i\binom{n+1}{i}(k-i)^n
       \qquad(1\le k\le n)
\]
and checks agreement with the recurrence. It then forms the matrix-square
coefficients
\[
 B(n,k)=\sum_{j=k}^n A(n,j)A(j,k),\qquad
 R_n(x)=P_n(x)/x.
\]
It checks $B(n,n)=1$, $B(n,1)=n!$, positivity of all coefficients of
$R_n$, the order and negativity of the certificate points, and
\[
 \operatorname{sgn}R_n(q_{n,i})=(-1)^i
       \qquad(1\le i\le n-1).
\]
Since $R_n(0)=n!>0$, these signs give one root in each of the $n-1$
disjoint intervals between consecutive listed points, starting at zero.
They exhaust the degree. Thus every root of $R_n$ is negative and simple;
$P_n$ has the additional simple root at zero. The case $P_0=1$ is immediate.

The sign evaluations clear denominators during Horner evaluation.
For $q=a/d$ with $d>0$ and
$R(x)=\sum_{i=0}^m r_i x^i$, the sign is the sign of the integer
\[
 d^mR(a/d)=\sum_{i=0}^m r_i a^i d^{m-i}.
\]
Starting with $h=r_m$, the verifier successively replaces
$h$ by $ah+r_i d^{m-i}$ for $i=m-1,\ldots,0$.
It retains the original rational Horner routine and checks equality of
the two evaluations, after multiplying by $d^m$, at every certificate
point in the first $20$ rows.
The verifier does not rely on the numerical procedure that originally
selected the rational points.

\subsection{Rational bounds for the analytic constants}

The scalar verifiers use positive Taylor series and geometric remainder
bounds. For rational $x\ge0$ and $M+2>x$, put
\[
 E_M(x)=\sum_{j=0}^{M}\frac{x^j}{j!}.
\]
Then
\[
 E_M(x)\le e^x\le E_M(x)+
 \frac{x^{M+1}}{(M+1)!}\frac1{1-x/(M+2)}.
\]
Taking reciprocals gives outward bounds for $e^{-x}$. To bound logarithms,
write $x=2^r y$ with $1\le y\le2$, and use
\[
 \log y=2\sum_{j=0}^{\infty}\frac{z^{2j+1}}{2j+1},
 \qquad z=\frac{y-1}{y+1}\in[0,1/3].
\]
After terms $0,\ldots,M-1$, the remaining sum is at most
\[
 \frac{2z^{2M+1}}{(2M+1)(1-z^2)}.
\]
The same formula bounds $\log2$, and reciprocation handles $x<1$.
The series
\[
 I_0(v)=\sum_{j=0}^{\infty}\frac{(v^2/4)^j}{(j!)^2}
\]
is treated in the same way, with successive tail ratios bounded by
$(v^2/4)/(M+2)^2$. Square-root comparisons are reduced to squared
rational inequalities with the signs checked first.

For inequalities on unbounded intervals, the proof first reduces the
claim to positivity of a polynomial in a shifted nonnegative variable;
the verifier checks its coefficients. On a bounded interval the
corresponding Bernstein coefficients are used. These are finite rational
calculations. The reduction from the original function to the polynomial
inequality is given in the relevant proof above.

Table~\ref{tab:scalar-summary} records the principal bounds used in the
proof. Every terminating decimal in the table denotes an exact rational
number. The estimates for $n\ge250$ are uniform in $n$; the program checks
the scalar reductions, rather than enumerating values of $n$.

\begin{table}[htbp]
\centering\small
\begin{tabular}{p{.46\textwidth}p{.43\textwidth}}
\toprule
Quantity and range & Certified bound \\
\midrule
Positive pole tail, $0\le s\le3\sqrt n$ & $<.85$ \\
Positive pole tail, $3\sqrt n\le s\le4\sqrt n$ & $<.33$ \\
Negative pole tail, $-4\sqrt n\le s\le0$ & $<.88$ \\
Variance parameters, $\pi\le u\le4\sqrt n$ &
 $Y>7.5$, $Y/n<.26$, $h=Y^2/n<17$ \\
Tilted variance & $\operatorname{Var}L<8/h$ \\
Quadratically weighted far tail in the variance bound & $<.003$ \\
Poisson--binomial maximum mass, $v\ge.7$ & $<.47/\sqrt v$ \\
Each weighted tail outside a paired rectangle & $<.04AB$ \\
Ideal mass outside an isolated mode & $<.907$ times the mode \\
Eulerian replacement error in the overlap band & $<AB/(1000H)$ \\
Independent-zero location error & $<.001$ in $\log X$ \\
Derivative near each independent zero & $>.128$ \\
Paired middle sign after all errors & $<-.79AB/H$ \\
Same-coordinate event spacing in the overlap band & $>12$ \\
\bottomrule
\end{tabular}
\caption{Principal rational constants in the proof.}
\label{tab:scalar-summary}
\end{table}

For example, write
\[
 S(a)=\sum_{j\ge1}j^2e^{-aj(j-1)},\qquad
 F(h)=.47\{S(h/7)+e^{-h/2}S(h/3)+.003\}.
\]
The maximum-mass argument for the variance reduces to
$F(h)<(8/h)^{3/2}$. On each interval $[a,b]$ in
Table~\ref{tab:variance-numerical}, monotonicity bounds $F(h)$ by the
listed rational upper bound for $F(a)$, whose square is less than
$(8/b)^3$. The range $0<h\le2$ is covered by
$S(a)<2a^{-3/2}$ for $a\le2/3$, as proved above.
For $11\le h<17$, the two-adjacent-point argument is checked on the
twelve intervals $[j/2,(j+1)/2]$, $j=22,\ldots,33$. On each interval,
the verifier checks the positive gap and the squared inequality used to
bound $\frac14+\Delta+\sqrt\Delta$ by $8/h$.

\begin{table}[htbp]
\centering\small
\begin{tabular}{ccc}
\toprule
$a$ & $b$ & Upper bound for $F(a)$ \\
\midrule
2&3&2.969\\
3&4&1.803\\
4&6&1.298\\
6&8&.861\\
8&10&.676\\
10&11&.584\\
\bottomrule
\end{tabular}
\caption{The six intermediate intervals in the variance estimate.}
\label{tab:variance-numerical}
\end{table}

For the edge bridge $250\le n\le1600$, put
$L=\lceil\sqrt n/4\rceil$. If $\ell\le L$, then $\ell\le10$ and
\[
 n\ge n_{\min}(\ell)=\max\{250,16(\ell-1)^2+1\}.
\]
The required upper bounds decrease with $n$ at fixed $\ell$, and the
affine lower bounds increase, so the
ten endpoint checks in Table~\ref{tab:finite-edge-endpoints} cover the
whole range. They give outside coordinate ratios
$(.036,.015,.0006,.015)$ in the four directions, total ideal error
$<.054$, and central Eulerian error $<10^{-6}$. Their sum is below the
quadratic sign margin $1/15$. The same checks give the phase comparison
and $40<\log X_L^-<4\sqrt n$.

\begin{table}[htbp]
\centering\small
\begin{tabular}{cccccccc}
\toprule
$\ell$&$1,2,3,4$&5&6&7&8&9&10\\
\midrule
$n_{\min}(\ell)$&250&257&401&577&785&1025&1297\\
\bottomrule
\end{tabular}
\caption{Endpoints used for the finite edge bridge.}
\label{tab:finite-edge-endpoints}
\end{table}

\subsection{Supplement map}

The reproduction entry point is \path{src/verify_all.py}. The required
files are organized as follows.
\begin{itemize}
\item \path{src/verify_certificates.py} and
\path{src/eulerian_core.py} reconstruct the Eulerian rows and check the
rational sign certificates.
\item Under \path{src/astra_literature/},
\path{verify_extended_variance_constants.py} checks the variance estimate,
\path{verify_joint_tail_audit_constants.py} checks both lattice tails and
the Eulerian replacement error, and \path{verify_finite_edge_bridge.py}
checks the finite-range edge and phase constants.
\item Under \path{src/astra_structural/},
\path{verify_positive_bulk_constants.py} and
\path{verify_positive_bulk_extension_constants.py} check the two positive
bulk ranges, while \path{verify_negative_bulk_pole_constants_4sqrt.py}
checks the negative bulk through $4\sqrt n$.
These scripts import their rational Taylor helpers from
\path{verify_negative_bulk_variance_constants.py}, which is also included.
The file \path{verify_paired_cluster_sign_constants.py} checks the
independent-zero and paired-sign constants.
\item Under \path{src/astra_interlacing/},
\path{verify_mode_geometry_constants.py} checks the overlap-band
geometry, outer separator, lattice-sum margins, and separator phase
constants.
\end{itemize}

The finite polynomial certificates and the scalar checks have distinct
roles. The former prove the initial rows directly. The latter verify the
explicit constants in the uniform analytic proof, whose root count
covers all remaining rows.

\bibliographystyle{amsplain}
\bibliography{references}
\end{document}
